\documentclass[11pt,a4paper]{article}
\usepackage[margin=25mm]{geometry}
\usepackage{amsmath,booktabs,siunitx,hyperref}
\title{CSA noise: capacitance and injection controls}
\author{NAUSICAA calculation record}
\date{\today}
\begin{document}\maketitle
The supplied 180 nm circuit is evaluated with the corrected Spice2 card,
including its NOIMOD=2 noise parameters. The external input capacitance is
varied; intrinsic MOS capacitances, transistor sizes and biases are retained.
The three external loads bracket the design assumptions requested by the author.
Topmetal-S reports approximately 3 pF for its total input capacitance.

\section*{Conditions}
The feedback capacitor is 3 fF and the temperature is 20 degrees Celsius.
Noise is stationary about the zero-signal operating point and integrated over
0.1 Hz--1 MHz. The ion current is the existing uniform mean arrival distribution
over 10 ms, for a 0.5 mm square pad in an infinite grounded plane and a finite
5 mm approach at 3 m/s. Its charge is $e\Delta\phi_w$, without renormalization.
The frame estimator subtracts centered 300 microsecond voltage averages at
the frame boundaries, with $\alpha=\exp(-10\,\mathrm{ms}/0.3\,\mathrm{s})$.
ENC is the filtered output RMS divided by the simulated one-ion signal gain.

\begin{center}\begin{tabular}{rr}\toprule
External capacitance (pF) & Frame ENC (electrons RMS)\\\midrule
0.75 & 232.3 \\
1.25 & 295.9 \\
3 & 521.7 \\
\bottomrule\end{tabular}\end{center}

\section*{Injection at 1.25 pF}
A pulse with a 1 microsecond plateau and 10 ns edges injects one elementary
charge at simulation time 5 ms. Frame windows are centered at 2 and 12 ms,
so collection occurs 3 ms into this frame.
The same 10 ms estimator gives an ENC of 292.1 electrons RMS, compared
with 295.9 for the uniform ion signal.

The trapezoidal filter is a difference of two moving averages of duration $T$,
separated by $2T$. Its rise and flat-top durations are both $T$.
The signal is evaluated at the noiseless-template peak time. Noise is
integrated with exactly the same filter; no pole-zero correction is applied.
At $T=300$ microseconds the ENC is 203.6 electrons RMS.

\begin{center}\begin{tabular}{rr}\toprule
Rise and flat top (microseconds) & Injection ENC (electrons RMS)\\\midrule
180 & 208.0 \\
200 & 207.1 \\
250 & 205.2 \\
300 & 203.6 \\
350 & 202.4 \\
\bottomrule\end{tabular}\end{center}

\section*{Noise source diagnosis}
At 1.25 pF, flicker contributes 96.04\% of the frame-noise variance.
The contributions of the individual transistors to total variance are:
\begin{center}\begin{tabular}{rr}\toprule
Transistor & Variance fraction (percent)\\\midrule
M1 & 0.21 \\
M2 & 2.32 \\
M3 & 11.83 \\
M4 & 81.77 \\
M5 & 0.00 \\
M6 & 0.24 \\
MF & 3.63 \\
\bottomrule\end{tabular}\end{center}
M4 has $W=0.72$ and $L=0.36$ micrometres in the supplied netlist.
This identifies a specific device for geometry and bias optimization.
The present calculation does not compare calibrated 130 and 180 nm process
cards, so it cannot attribute the difference from Topmetal-S to the process node.
The Topmetal-S measurement is approximately 115 electrons RMS after digital
trapezoidal filtering; its filter timing is not reported in that paper.

\section*{Validation and reproduction}
The analysis checks finite-path and injected charge, halves transient timesteps,
doubles the spectral density from 200 to 400 points per decade, increases
quadrature order from 8 to 16, and refines the injection peak-time grid.
The largest relative error among the frame checks is 0.000366.
Each injection convergence check is below $10^{-3}$ and its quadrature
check is below $10^{-5}$. Integrated transistor variances close to the total
within $10^{-3}$. The 1.25 pF frame result reproduces the preceding calculation
to a relative tolerance of $10^{-9}$.

Run \texttt{scripts/run\_spice2\_diagnostics.py} with \texttt{--simulator},
then \texttt{scripts/analyze\_spice2\_diagnostics.jl}, and
\texttt{scripts/present\_spice2\_diagnostics.py}, from the repository root.
Acquisition saves native decks and HDF5; Julia saves the numerical analysis;
Python generates this report from that analysis.

\begin{thebibliography}{9}
\bibitem{yang} Y. Yang et al., Design and preliminary test results of the
charge sensitive amplifier for gain-less charge readout in high-pressure TPC,
JINST 19 (2024) C03031, \url{https://doi.org/10.1088/1748-0221/19/03/C03031}.
\bibitem{bsim} W. Liu et al., BSIM3v3.2.2 MOSFET Model Users' Manual,
UC Berkeley (1999), Chapter 8.
\end{thebibliography}
\end{document}
